\subsubsection{Gestión de Análisis de PMS}

\noindent\textbf{Caso de uso:} Crear Repeticiones y Gestionar Tandas de PMS

\vspace{0.5em}
\noindent\textbf{Descripción:} El usuario crea repeticiones dentro de tandas, ingresa pesos, y el sistema calcula automáticamente promedios y coeficiente de variación, creando nuevas tandas si es necesario.

\vspace{0.5em}
\noindent\textbf{Pre-Condición:}
\begin{itemize}[nosep]
    \item El análisis de PMS debe existir con estado \texttt{EN\_PROCESO}
    \item La configuración inicial debe estar completa (repeticiones esperadas, tipo de semilla)
    \item Debe existir al menos una tanda
    \item El usuario debe estar autenticado
\end{itemize}

\vspace{0.5em}
\noindent\textbf{Post-Condición:}
\begin{itemize}[nosep]
    \item Repeticiones creadas con pesos registrados
    \item Cálculos automáticos realizados (promedio, CV)
    \item Si CV > umbral $\rightarrow$ nueva tanda creada automáticamente
    \item Análisis listo para finalizar cuando se cumplen condiciones
\end{itemize}

\vspace{0.5em}
\noindent\textbf{Actores:} Analista, Administrador.

\vspace{0.5em}
\noindent\textbf{Flujo de evento principal:}
\begin{enumerate}[nosep]
    \item \textbf{Fase 1: Acceso al análisis}
    \begin{itemize}[nosep]
        \item El usuario accede al detalle del análisis de PMS
        \item Ve la configuración establecida:
        \begin{itemize}[nosep]
            \item Número de repeticiones esperadas (4-20)
            \item Tipo de semilla (brozosa o normal)
            \item Umbral CV (6\% o 4\%)
        \end{itemize}
        \item Ve la tanda actual activa
    \end{itemize}
    \item \textbf{Fase 2: Creación de repeticiones}
    \begin{itemize}[nosep]
        \item El usuario selecciona ``Agregar repetición''
        \item Ingresa el peso en gramos (3 decimales)
        \item Guarda la repetición
        \item El sistema:
        \begin{itemize}[nosep]
            \item Crea la repetición
            \item La vincula a la tanda actual
            \item Incrementa el contador de repeticiones de la tanda
        \end{itemize}
    \end{itemize}
    \item \textbf{Fase 3: Cálculos automáticos}
    \begin{itemize}[nosep]
        \item Después de crear cada repetición, el sistema calcula:
        \begin{itemize}[nosep]
            \item Peso promedio de la tanda = suma de pesos / número de repeticiones
            \item Coeficiente de variación (CV) = (desviación estándar / promedio) $\times$ 100
        \end{itemize}
    \end{itemize}
    \item \textbf{Fase 4: Validación de CV y creación de tandas}
    \begin{itemize}[nosep]
        \item Si CV $\leq$ umbral:
        \begin{itemize}[nosep]
            \item La tanda queda válida
            \item Muestra indicador verde
        \end{itemize}
        \item Si CV > umbral:
        \begin{itemize}[nosep]
            \item El sistema crea automáticamente una nueva tanda
            \item Muestra mensaje: ``El CV excede el límite. Se creó automáticamente la Tanda [N]''
            \item La nueva tanda queda activa para agregar repeticiones
        \end{itemize}
        \item El proceso se repite hasta lograr CV $\leq$ umbral
    \end{itemize}
    \item \textbf{Fase 5: Finalización}
    \begin{itemize}[nosep]
        \item Cuando todas las repeticiones están completas y CV es válido:
        \begin{itemize}[nosep]
            \item El análisis queda listo para finalizar
            \item El usuario puede proceder a finalizar el análisis
        \end{itemize}
    \end{itemize}
\end{enumerate}

\vspace{0.5em}
\noindent\textbf{Flujos alternativos:}
\begin{itemize}[nosep]
    \item \textit{Intentar crear más repeticiones que las configuradas:}
    \begin{itemize}[nosep]
        \item El sistema valida el número de repeticiones esperadas
        \item Bloquea la creación si se excede
        \item Muestra mensaje: ``Ya se alcanzó el número máximo de repeticiones configuradas ([N])''
    \end{itemize}
    \item \textit{CV persistentemente alto:}
    \begin{itemize}[nosep]
        \item El sistema sigue creando tandas automáticamente
        \item El usuario puede revisar la configuración o datos ingresados
        \item Puede consultar con supervisor sobre la calidad de la muestra
    \end{itemize}
    \item \textit{Editar peso de repetición existente:}
    \begin{itemize}[nosep]
        \item El usuario puede editar repeticiones ya creadas
        \item El sistema recalcula automáticamente promedio y CV
        \item Puede generar nueva tanda si CV excede umbral
    \end{itemize}
    \item \textit{Eliminar repetición:}
    \begin{itemize}[nosep]
        \item El usuario puede eliminar repeticiones
        \item El sistema recalcula automáticamente
        \item Puede eliminar tandas vacías automáticamente
    \end{itemize}
\end{itemize}

\vspace{0.5em}
\noindent\textbf{Requerimientos especiales:}
\begin{itemize}[nosep]
    \item Cálculos automáticos de promedio y CV en tiempo real
    \item Creación automática de tandas cuando CV excede umbral
    \item No se puede modificar la configuración inicial (repeticiones esperadas, tipo de semilla)
    \item Validación de número máximo de repeticiones por tanda
    \item Pesos con 3 decimales de precisión
\end{itemize}
